\section{Literature review}\label{sec:literature}

\subsection{What makes a useful assessment question}\label{sec:literature-2-1}

A useful question assesses the intended knowledge or skill and gives learners enough information to answer it. Multiple-choice questions also need a defensible answer key and plausible wrong options. The NBME item-writing guide recommends focused questions with answer options that address the same task. It describes how ambiguous wording, irrelevant difficulty and unintended clues can weaken an item. Although written for health-sciences examinations, its guidance provides a starting point for reviewing generated questions \parencite{nbme2024}.

For program tracing, a plausible distractor can represent a mistake in following execution. \textcite{thomas2019} generate distractors from alternative execution paths, connecting the wrong answers to possible tracing errors instead of arbitrary values. Whether learners actually make those errors would need to be examined through their responses. The project can test whether distractors are wrong and distinct; their educational usefulness needs a different form of assessment.

Correctness is therefore one part of question quality. Clear wording and alignment with the intended skill matter even when execution confirms the answer. Difficulty also depends on the learners and conditions under which they attempt the question.

\subsection{Reusable templates and question variation}\label{sec:literature-2-2}

A reusable template defines the content shared by several questions and the rules for varying it. In program tracing, input values can change while the code structure stays fixed. This gives learners further practice with the same reasoning pattern without requiring a new question to be written each time.

PrairieLearn implements this approach through question files and generation functions. Its Python functions supply parameters and correct answers, and random generators are seeded from the question variant's seed \parencite{prairie-server}. Its AI authoring interface can generate those files from a natural-language request and revise them through follow-up instructions. Instructors are still asked to review the files, preview variants and submit test answers before use \parencite{prairie-ai}.

Template authoring concentrates effort at the start. \textcite{thomas2019} discuss the care required to identify suitable variations and the time needed to write and inspect templates. Reuse then spreads that effort across generated questions, but it also spreads any fault in the shared rules. A grader can repeatedly accept an incorrect reference answer. Checking a template should therefore examine how answers are computed, as well as whether its questions render and its grading code runs.

\subsection{Checking generated content and using feedback for revision}\label{sec:literature-2-3}

There are several ways to improve an unsuccessful generation. The model can produce a fresh candidate, review its earlier output, or revise it using evidence from external checks. Fresh generation offers another opportunity to succeed. Checking feedback adds information about what went wrong and, where possible, an input that exposes the fault.

Self-Refine uses the same language model to generate, critique and revise an output across several tasks \parencite{madaan2023}. Self-Debugging studies program correction using explanations and execution results. Its experiments found that unit-test feedback improved debugging performance where tests were available \parencite{chen2024}.

Property-based testing is particularly relevant to reusable rules. Developers describe properties that should hold and provide strategies for generating inputs. The test system searches for violations instead of relying only on individually written examples. Hypothesis provides this approach for Python and can simplify a failing generated input into a smaller counterexample \parencite{maciver2019}.

For an assessment template, a property might require the answer recipe to agree with execution for every tested input. Another might require all selected options to differ. A counterexample can make revision more specific than a general instruction to improve the question. Sampling keeps checking work bounded, while leaving some permitted inputs untested. A comparison with regeneration must also account for the extra candidate opportunities and computation used during revision.

\subsection{Methods for estimating question difficulty}\label{sec:literature-2-4}

Difficulty can be measured after learners answer a question or estimated beforehand. Classical item analysis includes the proportion answering correctly. Item response theory models responses in relation to learner ability and item characteristics. Both depend on response data, which may be unavailable for newly generated questions \parencite{nbme2024,brant2026}.

Before those data exist, authors can use a rubric based on the task's demands. Cognitive-load theory relates demand to elements that must be considered together, rather than simply the quantity of information \parencite{sweller2010}. For tracing, this suggests examining dependencies between variable updates and control flow, rather than counting lines alone. Turning that principle into easy, medium and hard labels still requires explicit rubric choices.

Tracing conditions affect performance too. \textcite{xie2018} taught novices to trace systematically and maintain an external record of variable values. Their participants performed better on the study's tracing tasks than the control group. An estimate should therefore state assumptions such as familiarity with the language and access to scratch paper.

Automated methods include directly asking a language model for a difficulty estimate and training a predictor using signals from models solving the questions. \textcite{zotos2025} evaluated both against student performance on machine-learning true/false questions. Their supervised approach used uncertainty signals from solving, illustrating a route that requires labelled training data rather than only a written rubric.

A model's own problem-solving ability is not a reliable substitute for evaluation of its difficulty estimates. \textcite{li2026} found that stronger problem solving did not consistently improve agreement with human difficulty. On Brazil's ENEM questions, \textcite{brant2026} found systematic underestimation alongside moderate ranking performance. The distinction is practical: answering a question and predicting how learners will struggle with it are different tasks.

Manually assigned references also need scrutiny. \textcite{hayes2007} explain why independent observer agreement matters when judgments are coded. Agreement with a project reference shows consistency with those labels; learner responses are needed to test whether the labels represent experienced difficulty.

\subsection{Methods for assessing learning-objective relevance}\label{sec:literature-2-5}

Relevance asks whether a question requires the skill specified by a learning objective. \textcite{zur2023} formulate question-to-objective labelling as a classification task and provide a benchmark drawn from science textbooks. In its OpenStax portion, questions inherit the objectives listed for their textbook subsection. This provides useful source data, but not individually verified evidence that each question requires every inherited objective.

Keyword matching offers a simple baseline. Embeddings represent meaning more flexibly and allow texts to be compared through cosine similarity. Sentence-BERT demonstrates an efficient approach in which the two representations can be computed independently \parencite{reimers2019}. Similarity alone leaves a gap between being about a topic and assessing its intended skill: a question containing sorting code may ask only for the list's length.

Joint scoring examines the question and objective together. Qwen3 rerankers take an instruction, query and document and use yes/no token scores to rank relevance \parencite{zhang2025}. The instruction can describe the intended relation, which makes this a candidate mechanism for objective matching. The Qwen3 work also supports instruction-conditioned embeddings, so the choice is not simply between an instructed model and embeddings that ignore instructions. Published retrieval benchmarks motivate these mechanisms; a paired educational evaluation is needed to choose between them for this project.

\subsection{Research gap and project positioning}\label{sec:literature-2-6}

Existing work provides reusable question generation and methods for revising model output. The remaining question for this project is how well a bounded authoring workflow can combine them: can checking feedback repair template faults, and does the resulting template provide useful reuse?

The project implements this workflow for Python tracing. Execution checks supply evidence about answers and distractors, while difficulty and relevance are evaluated against separate references. The evaluation distinguishes improvements in technical yield from the broader claim that generated questions are educationally better.
